\chapter{争议怎样改变两晋南北朝史的读法}
\label{app:controversies}

本附录不把所有疑问裁成标准答案。下列问题之所以重要，是因为若把后出的故事、整齐数字或现代概念直接倒投回二八〇至五八九年，读者会误解统一、迁徙、制度与再统一各自意味着什么。

\section{八王之乱与永嘉：不是一份精确的灾难统计}

\eventref{JH-0289-YANG-JUN-REGENT}{西晋·杨骏辅政}之后的政变、诸王相攻、洛阳陷落和长安失守，均有相当可靠的年代骨架。《晋书》与《资治通鉴》却把诏书、谋略、宫闱细节、伤亡和责任组织成高度连贯的兴亡叙事。可以确认中枢的继承与调兵失去约束；不可把每一封遗诏、每一位人物的内心或灾难数字写成现场档案。墓志、城址和地方材料可以校正局部时间和网络，仍不能还原全体军民经历。

\section{南渡与侨置：名目不等于人口普查}

\eventref{JH-0318-YUAN-EMPEROR}{东晋·建康称帝}确认晋的皇统在建康重立，侨州郡也能证明朝廷保留了迁徙记忆与行政安排。不能由此算出精确南渡人数，更不能把“侨姓”和“江东旧族”写作两个内部一致的群体。地方墓志、家族文本和城市遗址能展示一些家庭的迁居和仕宦，却保存的是可书写、可埋葬和被发现的少数；普通流民、军户和妇女的连续声音仍然稀少。

\section{十六国：国号、族名与疆界都在变化}

前赵、后赵、前秦、诸燕、后凉等名称多来自后出的整理，原始年号、称号和政权覆盖范围经常有异文。可以使用《晋书》载记、辑佚十六国史和《资治通鉴》建立相对次序；不能把它们画成同一时点边界精确的地图，也不能把羯、氐、羌、鲜卑等词转换成稳定的现代民族实体。城郭、部众、旧官吏、道路与地方资源的关系，往往比国号更能解释一地如何改换政治归属。

\section{北魏改革：诏令、执行与身份不是一回事}

\eventref{NSL-440-589-073}{北魏·均田令}、\eventref{NSL-440-589-075}{北魏·三长制}和\eventref{NSL-440-589-086}{北魏·迁都洛阳}可证明政策与战略转向的节点。它们不证明丈量、授田、基层组织、服饰、语言和婚姻在全境同步改变。所谓“汉化”若不分开宫廷规范、贵族策略、地方语言和家庭关系，只会抹平不同区域。墓志、造像和城址能显示部分群体如何参与，并不能替代乡村社会的全景。

\section{宗教政策：禁令不是信仰的总和}

太武帝灭佛与周武帝禁佛的诏令、寺院处置和僧尼还俗，都是可追踪的国家行动；地方执行范围、寺产去向和个人信仰变化必须另证。僧传、护教文字和后世传说尤其容易把冲突组织成善恶分明的故事。云冈、龙门、邺城造像和建康寺院证明宗教网络需要资源和赞助，却不能量化全部居民的虔诚或解释每次政策转向。

\section{五八九年：政治征服，不是区域差异消失}

\eventref{NSL-440-589-290}{隋·陈亡}是南北最高政权被重新接合的清晰节点。隋接收陈的州郡、户籍、仓储、官员、宗族和寺院，正说明统一之后仍须处理行政与社会的距离。陈亡不应被读作江南水路、地方家族、旧军镇和宗教资源从此同北方完全一致；它是另一段接收、调整和抵抗的开端。
